% Standalone copy-paste appendix fragment for Experiments 1.1--1.4.
% Requires the same table packages as main.tex, in particular booktabs and graphicx.
% This file intentionally inlines all table bodies and does not load external table files.

\subsection{Per-Seed Main Results}
\label{app:per-seed-main-results}

Because the main experiments use three seeds, we additionally report the
primary archive-quality and operational-assignment metrics at the
per-seed level. These appendix tables are intended to test whether the
main improvements are driven by a single favorable run or instead remain
directionally consistent across seeds.

\begin{table*}[t]
\centering
\small
\caption{Per-seed archive-quality breakdown for the three main seeds.
The table reports hypervolume and expected utility under the shared
four-objective evaluation protocol.}
\label{tab:app-per-seed-archive-quality}
\resizebox{\textwidth}{!}{%
\begin{tabular}{@{}llrr@{}}
\toprule
Seed & Method & Hypervolume ($\times 10^6$) & Expected Utility \\
\midrule
0007 & Constraint-Aware Stage-2 & 17.147 & -46.785 \\
0007 & Stage-1 Archive & 16.903 & -46.809 \\
0007 & Weighted-Sum & 6.705 & -136.944 \\
0007 & Preference-Conditioned PPO & 2.467 & -154.838 \\
0011 & Constraint-Aware Stage-2 & 27.255 & -36.658 \\
0011 & Stage-1 Archive & 25.894 & -42.718 \\
0011 & Weighted-Sum & 5.156 & -171.152 \\
0011 & Preference-Conditioned PPO & 5.314 & -175.946 \\
0019 & Constraint-Aware Stage-2 & 24.327 & -42.315 \\
0019 & Stage-1 Archive & 23.718 & -45.113 \\
0019 & Weighted-Sum & 5.654 & -176.605 \\
0019 & Preference-Conditioned PPO & 3.258 & -173.017 \\
\bottomrule
\end{tabular}
}
\end{table*}

\begin{table*}[t]
\centering
\small
\caption{Per-seed operational-assignment breakdown for the three main
seeds. The table reports feasible-assignment rate, mean violation, final
critical compromised hosts, and critical-impact count under the shared
acceptance protocol.}
\label{tab:app-per-seed-operational-assignment}
\resizebox{\textwidth}{!}{%
\begin{tabular}{@{}llrrrr@{}}
\toprule
Seed & Method & Feasible Rate & Violation & Final Critical Hosts & Critical Impacts \\
\midrule
0007 & Constraint-Aware Stage-2 & 0.025 & 0.278 & 0.000 & 0.000 \\
0007 & Stage-1 Archive & 0.000 & 0.321 & 0.000 & 0.000 \\
0007 & Weighted-Sum & 0.000 & 91.092 & 0.625 & 6.050 \\
0007 & Lagrangian PPO & 0.000 & 91.866 & 0.775 & 4.350 \\
0007 & Unconstrained Stage-2 & 0.350 & 1.816 & 0.000 & 0.000 \\
0011 & Constraint-Aware Stage-2 & 0.925 & 0.004 & 0.000 & 0.000 \\
0011 & Stage-1 Archive & 0.350 & 3.754 & 0.000 & 0.000 \\
0011 & Weighted-Sum & 0.000 & 102.630 & 0.875 & 4.925 \\
0011 & Lagrangian PPO & 0.000 & 103.041 & 0.900 & 1.000 \\
0011 & Unconstrained Stage-2 & 0.525 & 2.600 & 0.000 & 0.000 \\
0019 & Constraint-Aware Stage-2 & 0.975 & 0.000 & 0.000 & 0.000 \\
0019 & Stage-1 Archive & 0.575 & 1.856 & 0.000 & 0.000 \\
0019 & Weighted-Sum & 0.000 & 86.909 & 0.975 & 1.000 \\
0019 & Lagrangian PPO & 0.000 & 91.916 & 0.975 & 1.125 \\
0019 & Unconstrained Stage-2 & 0.500 & 1.967 & 0.000 & 0.000 \\
\bottomrule
\end{tabular}
}
\end{table*}

Across all three seeds, DA-CPSL remains the strongest or near-strongest
archive-quality method, showing that the main-table gains are not driven
by a single run. Per-seed deployment metrics likewise show that
DA-CPSL's feasible-assignment advantage and violation reduction persist
across seeds, rather than emerging from one unusually favorable seed.
The purpose of these tables is to confirm directional stability, not to
claim that all seeds produce identical magnitudes. The per-seed results
therefore strengthen the interpretation that the aggregate gains reflect
a repeatable trend under the stated protocol.

\subsection{Raw Pareto Assignment vs. Acceptable Pareto Assignment}
\label{app:raw-vs-acceptable-assignment}

\begin{table*}[t]
\centering
\small
\caption{Raw Pareto assignment versus acceptable Pareto assignment on the DA-CPSL Stage~2 archive. Feasible-assignment rate is measured over preference-grid assignments; violation and critical-risk audit metrics are assignment-weighted over selected policies.}
\label{tab:app-raw-vs-acceptable-assignment}
\resizebox{\textwidth}{!}{%
\begin{tabular}{@{}lrrrrr@{}}
\toprule
Assignment Rule & Expected Utility & Feasible Assign. Rate & Mean Violation & Final Critical Hosts & Critical Impact Count \\
\midrule
Raw Pareto Assignment & -41.919 & 0.679 & 1.189 & 0.000 & 0.000 \\
Acceptable Pareto Assignment & -43.706 & 1.000 & 0.000 & 0.000 & 0.000 \\
\bottomrule
\end{tabular}%
}
\end{table*}

This mechanism ablation uses the same DA-CPSL Stage~2 archive, the same
four-objective preference grid, and the same operational thresholds as
the main evaluation. The only change is the assignment rule: Raw Pareto
Assignment selects directly from the objective-space Pareto set, whereas
Acceptable Pareto Assignment first applies the operational acceptance
filter and then assigns from the acceptable Pareto set.

The raw rule preserves the highest preference utility, but it can assign
policies that are Pareto-optimal only in objective space while violating
the hard business or cost limits. In contrast, the acceptable rule
raises the feasible-assignment rate from \(0.679\) to \(1.000\) and
reduces mean violation from \(1.189\) to \(0.000\), at the cost of a
small reduction in expected utility. This appendix-only comparison does
not introduce a new training result and does not replace
Table~\ref{tab:operational-assignment-performance}; it isolates the
assignment-rule mechanism and directly supports the design motivation in
Section~\ref{subsec:acceptable-pareto-assignment}.

\subsection{Operational-Threshold Sensitivity}
\label{app:threshold-sensitivity}

\begin{table*}[t]
\centering
\small
\caption{Operational-threshold sensitivity under looser, default, and stricter business/cost limits. Feasible-assignment rate and mean violation are computed over preference-grid assignments from each method's raw Pareto candidates.}
\label{tab:app-threshold-sensitivity}
\resizebox{\textwidth}{!}{%
\begin{tabular}{@{}llrr@{}}
\toprule
Threshold Profile & Method & Feasible Assign. Rate & Mean Violation \\
\midrule
Looser & Constraint-Aware Stage-2 & 0.719 & 0.849 \\
Looser & Stage-1 Archive & 0.716 & 0.859 \\
Looser & Weighted-Sum & 0.000 & 101.242 \\
Looser & Lagrangian PPO & 0.000 & 95.034 \\
Looser & Unconstrained Stage-2 & 0.746 & 0.768 \\
\addlinespace
Default & Constraint-Aware Stage-2 & 0.679 & 1.189 \\
Default & Stage-1 Archive & 0.667 & 1.214 \\
Default & Weighted-Sum & 0.000 & 102.086 \\
Default & Lagrangian PPO & 0.000 & 95.825 \\
Default & Unconstrained Stage-2 & 0.746 & 1.069 \\
\addlinespace
Stricter & Constraint-Aware Stage-2 & 0.642 & 1.700 \\
Stricter & Stage-1 Archive & 0.304 & 2.217 \\
Stricter & Weighted-Sum & 0.000 & 104.273 \\
Stricter & Lagrangian PPO & 0.000 & 97.625 \\
Stricter & Unconstrained Stage-2 & 0.692 & 1.509 \\
\bottomrule
\end{tabular}%
}
\end{table*}

We also test whether the operational-assignment conclusion depends on a
single hand-picked threshold profile. Using the same official
four-objective artifacts, seeds, and preference grid, we evaluate three
business/cost profiles derived from the Table~B calibration set:
looser \(q=0.10\), default \(q=0.25\), and stricter \(q=0.50\); the
default profile exactly matches the shared thresholds used in the main
evaluation. For each profile, the assignment rule first selects from
each method's raw Pareto candidates by preference utility and then
audits whether the selected objective vector satisfies the corresponding
business and cost limits.

As expected, tighter thresholds reduce feasible-assignment rates and
increase mean violation for every method. DA-CPSL moves from
\((0.719, 0.849)\) under the looser profile to \((0.679, 1.189)\) at
the default profile and \((0.642, 1.700)\) under the stricter profile,
showing gradual degradation rather than a collapse at one operating
point. Stage-1 Only is more sensitive to tightening, dropping from
\(0.716\) feasible-assignment rate under the looser profile to
\(0.304\) under the stricter profile, while Unconstrained Stage~2
retains slightly higher feasibility but also accumulates more violation
as the limits tighten, from \((0.746, 0.768)\) to \((0.692, 1.509)\).
Weighted-Sum and Lagrangian PPO remain infeasible across all three
profiles, with mean violation above \(95\), so the appendix conclusion
is robustness of the relative deployment trend rather than a new
training or replay result. It isolates threshold sensitivity at the
candidate/objective assignment level.

\subsection{Zero-Observed Critical-Event Confidence Bounds}
\label{app:zero-event-confidence-bounds}

\begin{table*}[t]
\centering
\small
\caption{Zero-observed critical-event confidence bounds under the finite semantic replay protocol. Bounds are one-sided Clopper--Pearson upper bounds at 95\% confidence and do not imply zero risk.}
\label{tab:app-zero-event-confidence-bounds}
\resizebox{\textwidth}{!}{%
\begin{tabular}{@{}llrrr@{}}
\toprule
Event & Method & Observed Events / Episodes & Observed Rate & 95\% CP Upper Bound \\
\midrule
Ever Critical Breach & Constraint-Aware Stage-2 & 0/120 & 0.00\% & 2.47\% \\
Persistent Critical Breach & Constraint-Aware Stage-2 & 0/120 & 0.00\% & 2.47\% \\
Any Critical Impact & Constraint-Aware Stage-2 & 0/120 & 0.00\% & 2.47\% \\
Any Final Critical Host Compromised & Constraint-Aware Stage-2 & 0/120 & 0.00\% & 2.47\% \\
Ever Critical Breach & Stage-1 Archive & 0/120 & 0.00\% & 2.47\% \\
Persistent Critical Breach & Stage-1 Archive & 0/120 & 0.00\% & 2.47\% \\
Any Critical Impact & Stage-1 Archive & 0/120 & 0.00\% & 2.47\% \\
Any Final Critical Host Compromised & Stage-1 Archive & 0/120 & 0.00\% & 2.47\% \\
Ever Critical Breach & Unconstrained Stage-2 & 0/120 & 0.00\% & 2.47\% \\
Persistent Critical Breach & Unconstrained Stage-2 & 0/120 & 0.00\% & 2.47\% \\
Any Critical Impact & Unconstrained Stage-2 & 0/120 & 0.00\% & 2.47\% \\
Any Final Critical Host Compromised & Unconstrained Stage-2 & 0/120 & 0.00\% & 2.47\% \\
\bottomrule
\end{tabular}%
}
\end{table*}

The zero critical-event entries in the main audit are finite-protocol
observations, not formal guarantees of zero risk. We therefore report
one-sided Clopper--Pearson upper confidence bounds for event-style
semantic replay outcomes with zero observed events. For each method and
event, the analysis pools the three official seeds, giving
\(3\times 40=120\) semantic replay episodes under the same Table~B
protocol.

For DA-CPSL, zero observed ever-critical breaches, persistent critical
breaches, critical-impact events, and final critical compromised hosts
yield a \(95\%\) upper bound of \(2.47\%\) on the corresponding
per-episode event rate under this replay protocol. Stage-1 Only and
Unconstrained Stage~2 have the same zero-observed-event bound for these
audit events, whereas Weighted-Sum and Lagrangian PPO show nonzero
critical-event counts and are therefore retained only in the
machine-readable audit outputs for this analysis. This supports the
interpretation ``breach-free under the replay protocol, but not
risk-free.''
